\chapter{Unit 3: Collecting Data}

\section{Essential Sampling Methods \& Definitions}

\begin{table}[H]
\centering
\small
\begin{tabular}{@{}lp{4.2in}@{}}
\toprule
\textbf{Sampling Method} & \textbf{Key Procedure \& Purpose} \\
\midrule
\textbf{Simple Random Sample (SRS)} & Every group of $n$ individuals in the population has an equal chance of being selected. Eliminates selection bias. \\
\textbf{Stratified Random Sample} & Population divided into homogeneous groups (\emph{strata}) sharing similar traits; take an SRS from \emph{every} stratum. Reduces sampling variability. \\
\textbf{Cluster Sample} & Population divided into heterogeneous, naturally occurring groups (\emph{clusters}); randomly select several clusters and sample \emph{all} individuals in them. Saves time and cost. \\
\textbf{Systematic Sample} & Select every $k$-th individual from an ordered list starting from a randomly chosen initial point. \\
\textbf{Convenience / Voluntary} & Non-random sampling. Leads to severe bias. Never representative of the entire population. \\
\bottomrule
\end{tabular}
\caption{Summary of Sampling Methods}
\end{table}

\section{Observational Studies vs Experiments}

\begin{itemize}
    \item \textbf{Observational Study:} No treatments are imposed. Researchers observe or survey existing conditions. \emph{Can establish correlation/association only; NEVER causation.}
    \item \textbf{Experiment:} Treatments are actively imposed on experimental units. \emph{Can establish cause-and-effect conclusions if randomized assignment is used.}
\end{itemize}

\section{The 4 Principles of Experimental Design}

\begin{enumerate}
    \item \textbf{Comparison:} At least two treatment groups (or a treatment vs control group) to compare responses.
    \item \textbf{Random Assignment:} Experimental units are assigned to treatments at random to balance the effects of confounding variables across groups.
    \item \textbf{Replication:} Applying each treatment to enough experimental units so differences in effects can be distinguished from chance variation.
    \item \textbf{Control:} Keeping confounding variables constant across all treatment groups.
\end{enumerate}

\section{Blocking \& Matched Pairs Design}

\begin{itemize}
    \item \textbf{Randomized Block Design:} Experimental units are grouped into blocks based on an unavoidable confounding characteristic (e.g., age, pre-existing health condition) before treatments are randomly assigned within each block. \emph{"Block what you can control; randomize what you cannot."}
    \item \textbf{Matched Pairs Design:} A special case of blocking where subjects are paired by similarity and randomly assigned opposite treatments, or each subject receives both treatments in a randomized order (acting as their own control).
\end{itemize}

\begin{trapbox}[AP Exam Trap Alert: Confounding vs Bias]
\textbf{Confounding} occurs when an outside variable is associated with the explanatory variable AND affects the response variable, making it impossible to isolate the true effect of the explanatory variable. Always describe:
\begin{enumerate}
    \item How the confounding variable relates to the treatment/explanatory group.
    \item How it directly influences the response variable outcome.
\end{enumerate}
\end{trapbox}
